\chapter{Imagen solenoidal y resolución de Navier--Stokes}

\section{Dominio hidrodinámico y publicación Galerkin}

Se fija
\[
 H^s_\sigma(\mathbb T^3)=
 \{u\in H^s(\mathbb T^3;\mathbb R^3):
   \nabla\!\cdot u=0,\ \int_{\mathbb T^3}u=0\},
 \qquad s>\frac52,\quad\nu>0,
\]
con
$u_0\in C^\infty\cap H^s_\sigma$ y
$f\in C^\infty([0,\infty);H^\infty_\sigma)
\cap L^2_{\rm loc}([0,\infty);H^{s-1}_\sigma)$. Para la proyección
Leray--Galerkin $P_M$ se ponen
\[
 A_M=-P_M\Delta,
 \qquad B_M(v,w)=P_MP_\sigma((v\cdot\nabla)w),
 \qquad N_M(u):=B_M(u,u),
\]
y el campo finito satisface
\begin{equation}\label{eq:pdf2-ns-galerkin}
 \partial_tu_M+\nu A_Mu_M+B_M(u_M,u_M)=f_M,
 \qquad u_M(0)=P_Mu_0.
\end{equation}
La ecuación \eqref{eq:pdf2-ns-galerkin} es una publicación del estado, no su
definición completa.

Para que no cambie la notación a mitad de la cadena, se fijan una sola vez los
cuatro funcionales de energía de orden $s$:
\begin{equation}\label{eq:pdf2-ns-energy-functionals}
 \begin{aligned}
 \mathscr E_{s,M}(u)&:=\frac12\|\Lambda^su\|_{L^2}^2,
 &\mathscr D_{s,M}(u)&:=\|A_M^{1/2}\Lambda^su\|_{L^2}^2,\\
 \mathscr B_{s,M}(u)&:=
   \operatorname{Re}\langle\Lambda^sN_M(u),\Lambda^su\rangle,
 &\mathscr F_{s,M}(u,f)&:=
   \operatorname{Re}\langle\Lambda^sP_Mf,\Lambda^su\rangle.
 \end{aligned}
\end{equation}
El subíndice adicional $j$ denota los mismos cuatro funcionales evaluados en
la coordenada publicada del nivel de profundidad $j$; no introduce otros
funcionales.
La ecuación Galerkin implica exactamente
\begin{equation}\label{eq:pdf2-ns-energy-identity}
 \frac d{dt}\mathscr E_{s,M}(u_M)
 +\nu\mathscr D_{s,M}(u_M)+\mathscr B_{s,M}(u_M)
 =\mathscr F_{s,M}(u_M,f).
\end{equation}

\section{Objeto fuente y coordenadas que sí están construidas}

La imagen solenoidal no se añade a posteriori. Es exactamente la construcción
autónoma $\mathfrak G_{\rm sol}$ producida antes de este capítulo y publicada en
\eqref{eq:pdf2-g-sol-explicito}; conserva el estado
Galerkin, los mapas de refinamiento, los sectores orientados, el operador disipativo,
la advección proyectada, el carry, la counidad $E_9$, la inclusión $J_9$, las
proyecciones $P_9=J_9E_9$ y $Q_{\rm micro}=I-P_9$, la memoria, la ruta y el terminal.
En particular,
\[
 E_9J_9=I,
 \qquad P_9^2=P_9,
 \qquad Q_{\rm micro}=I-P_9.
\]
Estas identidades tipan la descomposición entre publicación macroscópica y memoria
microscópica; no constituyen por sí solas una cota global.

En una profundidad nonádica $j$, el estado conserva
\begin{equation}\label{eq:pdf2-ns-estado}
 \Xi_{M,j}^{\rm NS}=
 \bigl(u_M;q_{A,M,j},q_{V,M,j},T_{M,j};
 \begin{gathered}
 \text{tríadas, hoja, orientación, carry, frontera, ruta, archivo y memoria}
 \end{gathered}
 \bigr).
\end{equation}
El retorno enriquecido reúne todas las ramas antes de evaluar:
\begin{equation}\label{eq:pdf2-ns-retorno}
 \mathcal R_{M,J,t}^{\rm enr}
 =\mathcal R_{1,M}\mathfrak U_{M,J,t}^*,
 \qquad
 \mathcal R_{M,J,t}^{\rm enr}Z_{M,J,t}=u_M(t).
\end{equation}
Una rama aislada de norma $3^J$ no es el retorno.

Para $N>M$, el defecto de conmutación con la no linealidad es el carry
\begin{equation}\label{eq:pdf2-ns-carry}
 C_{M\leftarrow N}:=
 P_MB(u_N,u_N)-B_M(P_Mu_N,P_Mu_N).
\end{equation}
Ponerlo igual a cero sin probar conmutación elimina precisamente la interacción de
altas frecuencias que debe conservar el refinamiento.

Las dos memorias positivas se mantienen separadas. Con
$r=\pi_{\rm HMT}/\varphi_{\rm HMT}^2>1$,
\begin{equation}\label{eq:pdf2-ns-dh}
 D_M=(r-1)(\mathcal M_M^+-\mathcal M_M^-),
 \qquad
 H_M=(1+r)(\mathcal M_M^++\mathcal M_M^-),
\end{equation}
y
\[
 \dot D_M=\mathscr B_{s,M},
 \qquad
 \dot H_M=\frac{1+r}{r-1}|\mathscr B_{s,M}|.
\]
$D_M$ publica signo; $H_M$ conserva magnitud positiva. Fundirlos destruye la
orientación o la positividad.

\section{Torre PC--Fock completa y almacenamiento terminal}

La elevación no se trunca en grado dos:
\begin{equation}\label{eq:pdf2-ns-pcf}
 \mathcal V_M^{\rm PCF}(u)
 :=\bigoplus_{n\ge0}\frac{R_{\rm F}^n}{\sqrt{n!}}u^{\odot n},
 \qquad
 \|\mathcal V_M^{\rm PCF}(u)\|^2
 =e^{R_{\rm F}^2\|u\|_{H^s}^2}.
\end{equation}
El radio forma parte de la coordenada y no puede desaparecer del lector. Para
cualquier $\rho>0$ se fija
\begin{equation}\label{eq:pdf2-ns-normalized-degree-one-reader}
 \rho^{[\rho]}_1:=\rho^{-1}\pi_1:
 \Gamma_s^\rho(V)\longrightarrow V,
 \qquad
 \rho^{[\rho]}_1\mathcal V_\rho^{\rm PCF}(u)=u.
\end{equation}
El nivel inicial usa $\rho^{[R_{\rm F}]}_1$ y el nivel $j$ usa
$\rho^{[r_j]}_1$; esta normalización se mantiene al retirar la profundidad.
Sobre la diagonal coherente, $z_{n,M}=u_M^{\odot n}$ satisface
\begin{equation}\label{eq:pdf2-ns-carleman}
 \dot z_{n,M}=L_{n,M}z_{n,M}
 +\mathcal N_{n,n+1;M}z_{n+1,M}
 +\mathcal F_{n,n-1;M}(f_M)z_{n-1,M}.
\end{equation}
Los lectores de grados uno y dos son
\[
 z_M(u)=\nu^{1/2}A_M^{1/2}\Lambda^su,
 \qquad
 w_M(u)=\nu^{-1/2}A_M^{-1/2}\Lambda^sN_M(u),
\]
y Young da
\begin{equation}\label{eq:pdf2-ns-young-port}
 |\mathscr B_{s,M}(u)|\le
 \beta\nu\mathscr D_{s,M}(u)+\frac{\|w_M(u)\|^2}{4\beta}.
\end{equation}
La estimación uniforme de las filas de Fourier controla este lector cuadrático en
el estado inicial. No controla todavía sus propagaciones futuras.

Sea
\[
 \mathcal N_{M,j}:=
 \frac{r-1}{r+1}H_{M,j}-D_{M,j}^{\rm mem}
 =2(r-1)M_{M,j}^-\ge0.
\]
En un horizonte $T$, el almacenamiento terminal es
\begin{equation}\label{eq:pdf2-ns-storage}
 \begin{aligned}
 G_{M,j,T}(t)
 &=\mathscr E_{s,M,j}(t)+2(r-1)(M^-_{M,j}(T)-M^-_{M,j}(t)),\\
 \dot G_{M,j,T}+\nu\mathscr D_{s,M,j}+|\mathscr B_{s,M,j}|
 &=\mathscr F_{s,M,j}.
 \end{aligned}
\end{equation}
Para tipar el terminal se conserva en el estado solenoidal completo la coordenada
de hoja
\[
 q^-_{M,j}:=\sqrt{2(r-1)\dot M^-_{M,j}}
 \in L^2(0,T;\mathbb R).
\]
Sea $\mathscr X^{\rm sol}_{M,j,T}$ el espacio de historia enriquecida que contiene,
como sumandos ortogonales, la traza $2^{-1/2}\Lambda^su_{M,j}$ y $q^-_{M,j}$.
La columna terminal es el operador de coordenadas
\begin{equation}\label{eq:pdf2-ns-terminal-column}
 \begin{aligned}
 \mathcal E_{M,j,T}^{\rm term}(t):\mathscr X^{\rm sol}_{M,j,T}
 &\longrightarrow L^2_\sigma(\mathbb T^3)
   \oplus L^2(t,T;\mathbb R),\\
 \mathcal E_{M,j,T}^{\rm term}(t)\Xi
 &:=2^{-1/2}\Lambda^su_{M,j}(t)
 \oplus\mathbf1_{[t,T]}q^-_{M,j}.
 \end{aligned}
\end{equation}
Por construcción de las dos coordenadas ortogonales,
\begin{equation}\label{eq:pdf2-ns-terminal-column-gram}
 \|\mathcal E_{M,j,T}^{\rm term}(t)\Xi\|^2
 =G_{M,j,T}(t).
\end{equation}
El sumando de memoria futura es una salida terminal tipada, no una componente de la
raíz inicial; su coste se controlará sólo después de la telescopía de la columna
completa. Para una observación $\tau$, las filas de pérdida viven en
$[0,\tau]$ y esta memoria terminal en $(\tau,T]$; los soportes temporales son
disjuntos. En la estimación integral se toma $\tau=T$, y entonces el sumando
futuro es cero.

\section{Construcción source--first con lift mixto estado--fuerza}

La entrada se fija antes de todo corte, profundidad e historia:
\begin{equation}\label{eq:pdf2-ns-data-space}
 \mathcal D_T:=H^s_\sigma(\mathbb T^3)
 \times L^2(0,T;H^{s-1}_\sigma(\mathbb T^3)),
 \qquad d=(u_0,f).
\end{equation}
Como los campos tienen media nula, $A=-\Delta$ es positivo e invertible en el
sector solenoidal. Se introducen el esfuerzo, el flujo viscoso y las dos ondas
\begin{equation}\label{eq:pdf2-ns-force-waves}
 \begin{aligned}
 a_T(f)&:=\frac{1}{2\sqrt\nu}A^{-1/2}\Lambda^sP_\sigma f,
 &z_M(u)&:=\sqrt\nu A_M^{1/2}\Lambda^su,\\
 b_M(u,f)&:=z_M(u)-P_Ma_T(f),\\
 q_M^{B,+}(u)&:=\bigl(\mathscr B_{s,M}(u)\bigr)_+^{1/2},
 &q_M^{B,-}(u)&:=\bigl(-\mathscr B_{s,M}(u)\bigr)_+^{1/2},\\
 q_M^B(u)&:=\bigl(q_M^{B,+}(u),q_M^{B,-}(u)\bigr).
 \end{aligned}
\end{equation}
Todos pertenecen a $L^2$ sobre el intervalo correspondiente y
$\|q_M^B(u)\|_{\mathbb C^2}^2=|\mathscr B_{s,M}(u)|$.
La potencia cruzada está tipada por
\begin{equation}\label{eq:pdf2-ns-cross-power}
 \mathscr F_{s,M}=2\operatorname{Re}\langle z_M,P_Ma_T(f)\rangle,
 \qquad \nu\mathscr D_{s,M}=\|z_M\|^2.
\end{equation}
Por tanto \eqref{eq:pdf2-ns-storage}, y no una sustitución de la potencia por la
norma de la fuerza, da la identidad exacta de scattering
\begin{equation}\label{eq:pdf2-ns-scattering-balance}
 \dot G_{M,j,T}+\|b_M\|^2+\|q^B_M\|^2=\|P_Ma_T(f)\|^2.
\end{equation}

\subsection{Firma cronológica, suavizante uniforme y creación mixta}

El dato de fuerza vive en $F:=H^{s-1}_\sigma(\mathbb T^3)$, mientras que la
torre PC--Fock usa como espacio de una partícula
$V:=H^s_\sigma(\mathbb T^3)$. La identificación entre ambos no se deja a la
inclusión Galerkin, cuya norma crecería con $M$. En el sector de media nula se
fija el isomorfismo de orden menos uno
\begin{equation}\label{eq:pdf2-ns-uniform-smoother}
 \begin{aligned}
 \mathscr S&:=A^{-1/2}:F\longrightarrow V,\\
 \mathscr S_M&:=A_M^{-1/2}P_M=P_M\mathscr S,\\
 \sup_M\bigl(\|\mathscr S_M\|_{F\to V}
 +\|\mathscr S_M^{-1}\|_{P_MV\to P_MF}\bigr)&<\infty.
 \end{aligned}
\end{equation}
La última cota es la equivalencia uniforme de las normas de Sobolev en el
subespacio sin modo cero. Así, $g:=\mathscr Sf\in L^2(0,T;V)$ y no se usa la
inclusión no uniforme $P_MH^{s-1}\hookrightarrow P_MH^s$.

Sea $\Delta_k(T)=\{0<t_1<\cdots<t_k<T\}$. El espacio de firma de la fuerza
suavizada es
\begin{equation}\label{eq:pdf2-ns-signature-space}
 \mathfrak S_T(V):=\bigoplus_{k\geq0}L^2(\Delta_k(T);V^{\widehat\otimes k}),
 \qquad \Delta_0(T)=\{*\}.
\end{equation}
Para $g=\mathscr Sf\in L^2(0,T;V)$ se define, sin usar la solución,
\begin{equation}\label{eq:pdf2-ns-force-signature}
 \operatorname{Sig}_T(g)_0=1,
 \qquad
 \operatorname{Sig}_T(g)_k(t_1,\ldots,t_k)
 =g(t_k)\widehat\otimes\cdots\widehat\otimes g(t_1).
\end{equation}
La simetría del integrando escalar proporciona
\begin{equation}\label{eq:pdf2-ns-signature-norm}
 \|\operatorname{Sig}_T(g)\|^2
 =\sum_{k\geq0}\frac{\|g\|_{L^2(0,T;V)}^{2k}}{k!}
 =e^{\|g\|_{L^2(0,T;V)}^2}
 \le e^{c_{\mathscr S}^2\|f\|_{L^2(0,T;F)}^2}.
\end{equation}

El término forzado de \eqref{eq:pdf2-ns-carleman} no es un puerto aditivo en
$f$ solamente. Para $0<r_{\rm F}<R_{\rm F}$ se tipa mediante la creación mixta
\begin{equation}\label{eq:pdf2-ns-mixed-creation}
 \begin{aligned}
 \widetilde{\mathcal F}^{R_{\rm F},r_{\rm F}}_M:
 P_MV\widehat\otimes\Gamma_s^{R_{\rm F}}(P_MV)
 &\longrightarrow\Gamma_s^{r_{\rm F}}(P_MV),\\
 \widetilde{\mathcal F}^{R_{\rm F},r_{\rm F}}_M(g\otimes x)
 &:=a^\dagger(g)x.
 \end{aligned}
\end{equation}
En el grado $n$ su acción es precisamente
$g\otimes z_{n-1}\mapsto n g\odot z_{n-1}$. Así se conserva el factor mixto
$\mathscr Sf\otimes x$ que no puede producir una transición lineal sobre la suma directa
``estado $\oplus$ fuerza''.
Con las normas ponderadas de \eqref{eq:pdf2-ns-pcf},
\begin{equation}\label{eq:pdf2-ns-creation-uniform}
 \|\widetilde{\mathcal F}^{R,r}_M(g\otimes x)\|_{\Gamma^r}
 \le c(R,r)\|g\|_V\|x\|_{\Gamma^R},
 \qquad
 c(R,r):=\sup_{n\ge1}\sqrt n\,r\left(\frac rR\right)^{n-1}<\infty.
\end{equation}
Para iterar se fija $0<\theta<1$ y $r_j=R_{\rm F}\theta^j$; cada paso usa
$(R,r)=(r_j,r_{j+1})$. Entonces
$\sup_{j,M}c(r_j,r_{j+1})\le
R_{\rm F}\sup_{n\ge1}\sqrt n\,\theta^n<\infty$.

La torre se escribe en la coordenada conjugada
$v_M=\mathscr S_Mu_M$. Sea
$\mathscr S_M^{\Gamma}:=\bigoplus_{n\ge0}
\mathscr S_M^{\widehat\otimes n}$. El generador homogéneo conjugado es
\begin{equation}\label{eq:pdf2-ns-conjugated-generator}
 \mathcal A^{0,\mathscr S}_{M}
 :=\mathscr S_M^{\Gamma}\mathcal A^0_M
   (\mathscr S_M^{\Gamma})^{-1},
 \qquad
 (\mathcal A^{0,\mathscr S}_{M})_{n,m}
 =\mathscr S_M^{\widehat\otimes n}
  (\mathcal A^0_M)_{n,m}
  (\mathscr S_M^{-1})^{\widehat\otimes m}.
\end{equation}
Esta conjugación es source--first y no altera la ODE: el lector de grado uno
aplica $\mathscr S_M^{-1}$ al final.

Sea $\mathcal T_M^{0,\mathscr S}(t,s)$ el propagador de la torre conjugada, incluido
en la realización solenoidal completa. Para $R>r>0$ se define primero sobre
\begin{equation}\label{eq:pdf2-ns-radial-transfer}
 \delta_{R\to r}(x_0,x_1,x_2,\ldots)
 :=\bigl(x_0,(r/R)x_1,(r/R)^2x_2,\ldots\bigr),
 \qquad
 \delta_{R\to r}\mathcal V_R^{\rm PCF}(u)=\mathcal V_r^{\rm PCF}(u).
\end{equation}
Este operador es contractivo y fija el vacío. Sobre
\begin{equation}\label{eq:pdf2-ns-dyson-domain}
 \Gamma_s^R(P_MV)_{\rm alg}\widehat\otimes
 \mathfrak S_T(P_MV)_{\rm fin}
\end{equation}
Sea $\mathfrak S_{T,{\rm fin}}^{\rm dec}$ el span de los núcleos elementales
\[
 g_k(t_1,\ldots,t_k)=
 g_k^{(k)}(t_k)\widehat\otimes\cdots\widehat\otimes g_k^{(1)}(t_1).
\]
Es denso en cada sumando $L^2(\Delta_k(T);V^{\widehat\otimes k})$. En la
fórmula siguiente los símbolos $g_k^{(i)}$ se leen en cada núcleo elemental y
la acción se extiende linealmente a su span; por tanto no se factoriza
arbitrariamente un tensor entrelazado. El lector cronológico toma valores en
$\Gamma_s^r(P_MV)$ y se define por
\begin{equation}\label{eq:pdf2-ns-dyson-reader}
 \begin{aligned}
 \mathscr D_{M,k}^{R\to r}(t)(x\otimes g_k)
 :={}&\delta_{R\to r}\int_{\Delta_k(t)}
 \mathcal T_M^{0,\mathscr S}(t,t_k)
 \widetilde{\mathcal F}_M(P_Mg_k^{(k)}(t_k)\otimes\cdot)\cdots\\
 &\hspace{18mm}\cdots\mathcal T_M^{0,\mathscr S}(t_2,t_1)
 \widetilde{\mathcal F}_M(P_Mg_k^{(1)}(t_1)\otimes\cdot)
 \mathcal T_M^{0,\mathscr S}(t_1,0)x\,d\mathbf t,\\
 \mathscr D_{M,T}^{R\to r}(t)&:=\sum_{k\geq0}
 \mathscr D_{M,k}^{R\to r}(t).
 \end{aligned}
\end{equation}
Al distribuir la pérdida total $R-r$ entre las $k$ creaciones y los propagadores,
la estimación analítica de escala de
\eqref{eq:pdf2-ns-creation-uniform} da, para constantes finitas dependientes del
corte pero no del núcleo elemental,
\begin{equation}\label{eq:pdf2-ns-dyson-homogeneous-bound}
 \|\mathscr D_{M,k}^{R\to r}(t)(x\otimes g_k)\|_{\Gamma^r}
 \le A_{M,R,r,T}\,
 \frac{B_{M,R,r,T}^{k}}{\sqrt{\Gamma(k+1)}}
 \|x\|_{\Gamma^R}\|g_k\|_{L^2(\Delta_k;V^{\widehat\otimes k})}.
\end{equation}
El factor $1/\Gamma(k+1)$ del volumen del simplex y Cauchy--Schwarz en la
suma de grados dan
\begin{equation}\label{eq:pdf2-ns-dyson-bound}
 \|\mathscr D_{M,T}^{R\to r}(t)(x\otimes\sigma)\|_{\Gamma^r}
 \le C_{M,R,r,T}\|x\|_{\Gamma^R}\|\sigma\|_{\mathfrak S_T},
 \qquad C_{M,R,r,T}<\infty.
\end{equation}
Por ello la serie converge en el codominio indicado y determina por densidad
un operador acotado para cada corte. No se usa aquí una cota uniforme en $M$:
esa uniformidad se obtendrá después de la identidad local de la coligación.
Al evaluarlo en
$x\otimes\operatorname{Sig}_T(\mathscr Sf)$ se obtiene la serie de Dyson del generador no
autónomo
\begin{equation}\label{eq:pdf2-ns-nonautonomous-generator}
 \mathcal A_M^{f,\mathscr S}(t)
 =\mathcal A_M^{0,\mathscr S}+a^\dagger(\mathscr S_MP_Mf(t)),
 \qquad
 \partial_tU_M^{f,\mathscr S}(t,s)
 =\mathcal A_M^{f,\mathscr S}(t)U_M^{f,\mathscr S}(t,s).
\end{equation}
La diferenciación de \eqref{eq:pdf2-ns-dyson-reader} en el core algebraico da
la conjugada de \eqref{eq:pdf2-ns-carleman}; la unicidad de la ODE finita da
\begin{equation}\label{eq:pdf2-ns-dyson-degree-one}
 \mathscr S_M^{-1}\rho^{[r]}_1
 \mathscr D_{M,T}^{R_{\rm F}\to r}(t)
 \bigl(\mathcal V_M^{\rm PCF}(\mathscr S_MP_Mu_0)
       \otimes\operatorname{Sig}_T(\mathscr S_MP_Mf)\bigr)
 =u_M(t).
\end{equation}
La igualdad incluye $t=0$: el factor $r^{-1}$ de
\eqref{eq:pdf2-ns-normalized-degree-one-reader} cancela exactamente la
coordenada radial $r\mathscr S_MP_Mu_0$ producida por
\eqref{eq:pdf2-ns-radial-transfer}. No se impone $R_{\rm F}=1$.
Esta construcción precede a la historia y no invoca
$\mathsf S_{M,T}d$ para definirla.

Sea
$\mathfrak s_{\rm sol}:H^s_\sigma\to\mathcal F_{\rm seed}^{\rm sol}$ la
restricción solenoidal del mapa de semilla: conserva hoja, orientación, residuo,
cociente y ruta basada y depende sólo de $u_0$. La raíz gobernante no se
redefine aquí: es exactamente la raíz $J_T^{\rm sol}$ de
\eqref{eq:pdf2-g-sol-raiz-dato}. Fijamos la identidad de nombres
\begin{equation}\label{eq:pdf2-ns-root-identification}
 J_T^{\rm mix}:=J_T^{\rm sol}:
 \mathcal D_T=\mathcal D_T^{\rm NS}\longrightarrow\mathcal H_T^{\rm sol}.
\end{equation}
Su publicación de coordenadas mixtas, y no una segunda raíz, es
\begin{equation}\label{eq:pdf2-ns-root-map}
 \begin{aligned}
 \operatorname{Coord}_T^{\rm mix}(u_0,f)&:=
 \bigl[\mathcal V^{\rm PCF}(u_0)
 \widehat\otimes\operatorname{Sig}_T(\mathscr Sf)\bigr]
 \oplus\mathfrak s_{\rm sol}(u_0),\\
 \operatorname{Coord}_T^{\rm mix}&:\mathcal D_T\longrightarrow
 \mathfrak H_T^{\rm coord},\\
 \mathfrak H_T^{\rm coord}&:=
 \bigl[\Gamma_s^{R_{\rm F}}(V)
 \widehat\otimes\mathfrak S_T(V)\bigr]
 \oplus\mathcal F_{\rm seed}^{\rm sol}.
 \end{aligned}
\end{equation}
La carta es una proyección de coordenadas del estado aceptado. En el subespacio
cíclico generado por la raíz se escribe
\begin{equation}\label{eq:pdf2-ns-mixed-coordinate-chart}
 \chi_T^{\rm mix}:\mathcal H_T^{\rm sol}\longrightarrow
 \mathfrak H_T^{\rm coord},
 \qquad
 \chi_T^{\rm mix}J_T^{\rm sol}(u_0,f)
 =\operatorname{Coord}_T^{\rm mix}(u_0,f).
\end{equation}
Su acción fuera del subespacio cíclico es la proyección ortogonal de
$\mathfrak G_{\rm sol}$ sobre la coordenada torre--firma--semilla. No se crea
un cociente nuevo ni se identifica la norma de esta carta con la norma total
de $\mathcal H_T^{\rm sol}$.
Con la normalización de la semilla,
\begin{equation}\label{eq:pdf2-ns-root-norm}
 \|\operatorname{Coord}_T^{\rm mix}(u_0,f)\|^2
 \leq e^{R_{\rm F}^2\|u_0\|_{H^s}^2
             +c_{\mathscr S}^2\|f\|_{L^2(0,T;H^{s-1})}^2}
 +c_{\rm seed}(1+\|u_0\|_{H^s}^2).
\end{equation}
Esta estimación prueba que la carta de coordenadas se forma sólo con el dato;
no transporta una norma nueva a $\mathcal H_T^{\rm sol}$ ni se usa para fabricar
la cota gobernante. La norma y la columna de $J_T^{\rm mix}=J_T^{\rm sol}$ son
las ya fijadas por $\mathfrak G_{\rm sol}$.

\subsection{Dato único de partida y núcleos densos}

El único dato admitido en este capítulo es la imagen completa
$\mathfrak G_{\rm sol}$ reproducida en
\eqref{eq:pdf2-g-sol-explicito}--\eqref{eq:pdf2-g-sol-cota-dato}. En particular,
se usan sin renombrarlos la raíz $J_T^{\rm sol}$, el codificador
$\Lambda_{M,T}$, las transiciones $\Gamma_{M,j}^{\rm term}$, la fila completa
$L_{M,j,T}^{\rm term}$, el terminal $\mathcal E_{M,j,T}^{\rm term}$, la columna
$\mathbf G_{M,j;J,T}^{\rm term}$ y la isometría $I_{M,J}$. No se introduce una
segunda realización de esos objetos.

Sea $\mathcal D_{T,{\rm alg}}\subset\mathcal D_T$ el subespacio de datos con
$u_0$ trigonométrico y $f$ suave, de soporte temporal compacto y con valores
trigonométricos. Es denso en $\mathcal D_T$. Los núcleos sobre los que se
escriben todas las acciones son
\begin{equation}\label{eq:pdf2-ns-cyclic-cores}
 \begin{aligned}
 \mathscr C_T^{\rm src}
 &:=\operatorname{span}\{J_T^{\rm sol}d:d\in\mathcal D_{T,{\rm alg}}\}
 \subset\mathcal H_T^{\rm sol},\\
 \mathscr C_{M,j,T}^{\rm sol}
 &:=\operatorname{span}\{x_{M,j,T}(d):d\in\mathcal D_{T,{\rm alg}}\},\\
 x_{M,0,T}(d)&:=\Lambda_{M,T}J_T^{\rm sol}d,\\
 x_{M,j+1,T}(d)&:=\Gamma_{M,j}^{\rm term}x_{M,j,T}(d).
 \end{aligned}
\end{equation}
Los espacios $\mathscr X_{M,j,T}^{\rm sol}$ son las completaciones de estos
núcleos en las normas ya contenidas en $\mathfrak G_{\rm sol}$. En particular,
ninguna norma se define restando la pérdida que después se quiere acotar.

La identificación exacta entre la notación de aplicación y la del dato es
\begin{equation}\label{eq:pdf2-ns-exact-datum-identification}
 \boxed{\begin{gathered}
 J_T^{\rm mix}=J_T^{\rm sol},\qquad
 \Gamma_{M,j,T}=\Gamma_{M,j}^{\rm term},\qquad
 L_{M,j,T}=L_{M,j,T}^{\rm term},\\
 \mathcal G_{M,J,T}^{\rm sol}=\mathbf G_{M,0;J,T}^{\rm term}
 \end{gathered}}
\end{equation}
Esta igualdad es literal, no una semejanza de nombres. La acción source--first
queda por tanto
\begin{equation}\label{eq:pdf2-ns-source-encoder-action}
 \boxed{\begin{gathered}
 d\xmapsto{J_T^{\rm sol}}J_T^{\rm sol}d
 \xmapsto{\Lambda_{M,T}}x_{M,0,T}(d)
 \xmapsto{\Gamma_{M,0:j,T}}x_{M,j,T}(d)
 \end{gathered}},
 \qquad
 \Gamma_{M,0:j,T}:=\Gamma_{M,j-1,T}\cdots\Gamma_{M,0,T}.
\end{equation}
Ninguna de estas tres flechas recibe $u_M$, una pérdida o un terminal como
entrada. La carta PC--Fock--firma de
\eqref{eq:pdf2-ns-root-map} publica las coordenadas de la primera flecha; la
igualdad de grado uno
\eqref{eq:pdf2-ns-dyson-degree-one} prueba, en el núcleo y luego por densidad,
\begin{equation}\label{eq:pdf2-ns-sol-transition-action}
 \pi_{M,j,T}^{(1)}x_{M,j,T}(u_0,f)=u_M(t_j),
 \qquad
 \pi_{M,j,T}^{(1)}:=
 \mathscr S_M^{-1}
 (\rho_1^{[r_j]}\widehat\otimes\varepsilon_{[t_j,T]})\pi_{M,j,T}^{\rm tf}.
\end{equation}
El factor $r_j^{-1}$ se aplica antes de $\mathscr S_M^{-1}$ y cancela el radio
de grado uno. Por ello el lector devuelve $u_M$, no $r_ju_M$ ni
$\mathscr S_Mu_M$.

Para $N\ge M$, la proyección de refinamiento actúa tensorialmente sobre la
carta de coordenadas:
\begin{equation}\label{eq:pdf2-ns-refinement-action}
 \begin{aligned}
 \mathfrak p_{M\leftarrow N}^{\Gamma}
 (v_1\widehat\otimes\cdots\widehat\otimes v_k)
 &:=(P_Mv_1)\widehat\otimes\cdots\widehat\otimes(P_Mv_k),\\
 \mathfrak p_{M\leftarrow N}^{\mathfrak S}
 (g_1\widehat\otimes\cdots\widehat\otimes g_k)
 &:=(P_Mg_1)\widehat\otimes\cdots\widehat\otimes(P_Mg_k).
 \end{aligned}
\end{equation}
Junto con la proyección de semilla de $\mathfrak G_{\rm sol}$, estas fórmulas
dan el mapa $\mathfrak p_{M\leftarrow N}^{\rm sol}$ y los cuadrados
\begin{equation}\label{eq:pdf2-ns-refinement-commutation}
 \begin{aligned}
 \mathfrak p_{M\leftarrow N}^{\rm sol}\Lambda_{N,T}J_T^{\rm sol}
 &=\Lambda_{M,T}J_T^{\rm sol},\\
 \mathfrak p_{M\leftarrow N}^{\rm sol}\Gamma_{N,j,T}
 -\Gamma_{M,j,T}\mathfrak p_{M\leftarrow N}^{\rm sol}
 &=\iota_{\rm carry}\mathcal C_{M\leftarrow N,j},\\
 \mathcal C_{M\leftarrow N,j}x_{N,j,T}(d)
 &=C_{M\leftarrow N}(d)|_{I_j}.
 \end{aligned}
\end{equation}
El miembro derecho es el operador de coordenada cuyo valor publicado es
exactamente el carry de
\eqref{eq:pdf2-ns-carry}; no se lo declara nulo para fabricar naturalidad.

\subsection{Las seis filas como componentes del único operador de pérdida}

Fijemos $0=t_0<t_1<\cdots<t_J=\tau\le T$, $I_j=[t_j,t_{j+1}]$ y el calendario
\begin{equation}\label{eq:pdf2-ns-cutoff-calendar}
 M_j:=3^jM,
 \qquad P_{M_j}\longrightarrow I
 \quad\hbox{fuertemente en cada }H^r_\sigma.
\end{equation}
Los seis codominios, con sus normas propias, son
\begin{equation}\label{eq:pdf2-ns-six-loss-spaces}
 \begin{aligned}
 \mathcal Y_{M,j}^{\rm diss}&=L^2(I_j;L^2_\sigma),&
 \mathcal Y_{M,j}^{\rm adv}&=L^2(I_j;\mathbb C^2),\\
 \mathcal Y_{M,j}^{\rm carry}&=L^2(I_j;H^{s-1}_\sigma),&
 \mathcal Y_{M,j}^{\rm micro}&=L^2(I_j;\operatorname{Ran}Q_{\rm micro}),\\
 \mathcal Y_{M,j}^{\rm shell}
 &=L^2(I_j;\operatorname{Ran}(P_{M_{j+1}}-P_{M_j})\cap H^{s-1}_\sigma),&
 \mathcal Y_{M,j}^{\rm mem}
 &=L^2(I_j;\operatorname{Ran}Q^{\rm mem}_{M,j+1}).
 \end{aligned}
\end{equation}
Se pone
$\mathcal Y_{M,j}^{\rm sol}:=\widehat\bigoplus_{\mathfrak a}
\mathcal Y_{M,j}^{\mathfrak a}$, donde
$\mathfrak a\in\{{\rm diss,adv,carry,micro,shell,mem}\}$, y se denota por
$P_{\mathfrak a}$ su proyección ortogonal. La fila aceptada de
$\mathfrak G_{\rm sol}$ se descompone, y no se duplica, mediante
\begin{equation}\label{eq:pdf2-ns-loss-parseval}
 L_{M,j,T}^{\rm term}:
 \mathscr X_{M,j,T}^{\rm sol}\longrightarrow\mathcal Y_{M,j}^{\rm sol},
 \qquad
 L_{M,j,T}^{\mathfrak a}:=P_{\mathfrak a}L_{M,j,T}^{\rm term}.
\end{equation}

Para escribir las acciones sin convertir funciones no lineales de $u$ en
operadores lineales ficticios, se usan las proyecciones de coordenadas del lift
mixto ya contenido en $\mathfrak G_{\rm sol}$:
\begin{equation}\label{eq:pdf2-ns-row-coordinate-readers}
 \begin{aligned}
 \pi_zx_{M,j,T}(d)&=\mathbf1_{I_j}\sqrt\nu A_M^{1/2}\Lambda^su_M,&
 \pi_ax_{M,j,T}(d)&=\mathbf1_{I_j}P_Ma_T(f),\\
 \pi_{B,\pm}x_{M,j,T}(d)&=\mathbf1_{I_j}q_M^{B,\pm}(u_M),&
 \pi_{\rm carry}x_{M,j,T}(d)&=\mathbf1_{I_j}C_{M_j\leftarrow M_{j+1}},\\
 \pi_{\rm micro}x_{M,j,T}(d)&=\mathbf1_{I_j}Q_{\rm micro}y^{\rm micro}_{M,j},&
 \pi_{\rm shell}x_{M,j,T}(d)&=\mathbf1_{I_j}
 (P_{M_{j+1}}-P_{M_j})N_{M_{j+1}}(u_{M_{j+1}}),\\
 \pi_{\rm mem}x_{M,j,T}(d)&=\mathbf1_{I_j}
 Q^{\rm mem}_{M,j+1}m^{\rm raw}_{M,j+1}.&&
 \end{aligned}
\end{equation}
Aquí $y^{\rm micro}_{M,j}$ y $m^{\rm raw}_{M,j+1}$ no son símbolos
externos: son respectivamente las proyecciones sobre las coordenadas
$Q_{\rm micro}$ y $\operatorname{Mem}_{M,j+1}$ del estado
\eqref{eq:pdf2-g-sol-explicito}. Las ocho aplicaciones $\pi$ son lineales en
el estado levantado, aunque sus publicaciones sobre la diagonal física sean
polinomiales o contengan una raíz positiva. Su dominio es el núcleo
\eqref{eq:pdf2-ns-cyclic-cores}; su extensión acotada es la componente
correspondiente de $L_{M,j,T}^{\rm term}$.

Con estas proyecciones, las seis acciones son materialmente
\begin{equation}\label{eq:pdf2-ns-six-loss-actions}
 \begin{aligned}
 L^{\rm diss}_{M,j,T}&=\pi_z-\pi_a,&
 L^{\rm adv}_{M,j,T}&=(\pi_{B,+},\pi_{B,-}),\\
 L^{\rm carry}_{M,j,T}&=\pi_{\rm carry},&
 L^{\rm micro}_{M,j,T}&=\pi_{\rm micro},\\
 L^{\rm shell}_{M,j,T}&=\pi_{\rm shell},&
 L^{\rm mem}_{M,j,T}&=\pi_{\rm mem}.
 \end{aligned}
\end{equation}
La acotación no se infiere de los nombres. Para
$x\in\mathscr C_{M,j,T}^{\rm sol}$, la columna aceptada da
\begin{equation}\label{eq:pdf2-ns-row-boundedness}
 \|L_{M,j,T}^{\mathfrak a}x\|^2
 \le\|L_{M,j,T}^{\rm term}x\|^2
 \le\|\mathbf G_{M,j;J,T}^{\rm term}x\|^2.
\end{equation}
Por densidad, cada fila tiene una única extensión acotada a
$\mathscr X_{M,j,T}^{\rm sol}$. Ya no se forma un cociente nuevo ni se define
la norma del dominio desde la fila.

La ausencia de sobreconteo es la identidad de Parseval del único operador
$L^{\rm term}$:
\begin{equation}\label{eq:pdf2-ns-six-loss-no-overcount}
 \boxed{\begin{gathered}
 \|L_{M,j,T}^{\rm term}x\|^2
 =\sum_{\mathfrak a}\|L_{M,j,T}^{\mathfrak a}x\|^2,
 \qquad x\in\mathscr X_{M,j,T}^{\rm sol}
 \end{gathered}}
\end{equation}
Esta ortogonalidad coincide con la geometría física: los intervalos $I_j$
son disjuntos; carry vive bajo $P_{M_j}$ y shell bajo
$P_{M_{j+1}}-P_{M_j}$; micro vive bajo $Q_{\rm micro}$; y
$Q^{\rm mem}_{M,j+1}$ elimina la memoria heredada y la copia del canal
advectivo. Además $q_M^{B,+}q_M^{B,-}=0$ puntualmente y
$\|(q_M^{B,+},q_M^{B,-})\|^2=|\mathscr B_{s,M}|$, de modo que las dos hojas
son un solo canal orientado.

\subsection{Terminal, columna y Stein derivados sin redefinir normas}

En el núcleo terminal, la acción del operador aceptado es
\begin{equation}\label{eq:pdf2-ns-terminal-action-on-core}
 \mathcal E_{M,J,T}^{\rm term}(\tau)x_{M,J,T}(d)
 =2^{-1/2}\Lambda^su_M(\tau)
 \oplus\mathbf1_{[\tau,T]}q^-_{M,J}.
\end{equation}
Su extensión acotada a $\mathscr X_{M,J,T}^{\rm sol}$ es parte de
$\mathfrak G_{\rm sol}$ y su Gram terminal es, sin complementos inventados,
\begin{equation}\label{eq:pdf2-ns-terminal-metric}
 U_{M,J,T}:=
 (\mathcal E_{M,J,T}^{\rm term})^*\mathcal E_{M,J,T}^{\rm term},
 \qquad
 \|\mathcal E_{M,J,T}^{\rm term}x_{M,J,T}(d)\|^2
 =G_{M,J,T}(\tau).
\end{equation}
El lector de la primera coordenada terminal está definido en todo su
codominio por
\begin{equation}\label{eq:pdf2-ns-terminal-reader}
 \mathcal R_{M,\tau}^{(1)}(h\oplus q):=
 \sqrt2\,\Lambda^{-s}h,
 \qquad
 \mathcal R_{M,\tau}^{(1)}
 \mathcal E_{M,J,T}^{\rm term}(\tau)x_{M,J,T}(d)=u_M(\tau).
\end{equation}
No se invierte $\Gamma$, no se elige representante de un cociente y no aparece
un operador terminal complementario sin acción.

La columna terminal se fija primero por
\begin{equation}\label{eq:pdf2-ns-terminal-tail-column}
 \mathbf G_{M,J;J,T}^{\rm term}:=\mathcal E_{M,J,T}^{\rm term}.
\end{equation}
Para $0\le j<J$, la columna de cola se escribe con todos sus renglones, en el
orden terminal y luego profundidad decreciente:
\begin{equation}\label{eq:pdf2-ns-complete-column}
 \mathbf G_{M,j;J,T}^{\rm term}:=
 \begin{bmatrix}
  \mathcal E_{M,J,T}^{\rm term}\Gamma_{M,j:J,T}\\
  L_{M,J-1,T}^{\rm term}\Gamma_{M,j:J-1,T}\\
  \vdots\\
  L_{M,j+1,T}^{\rm term}\Gamma_{M,j:j+1,T}\\
  L_{M,j,T}^{\rm term}
 \end{bmatrix},
 \qquad
 \Gamma_{M,j:k,T}:=\Gamma_{M,k-1,T}\cdots\Gamma_{M,j,T}.
\end{equation}
Para $j<J$ hay una igualdad de operadores entre columnas ya construidas,
\begin{equation}\label{eq:pdf2-ns-tail-column-recursion}
 \mathbf G_{M,j;J,T}^{\rm term}
 =\begin{bmatrix}
   \mathbf G_{M,j+1;J,T}^{\rm term}\Gamma_{M,j,T}\\
   L_{M,j,T}^{\rm term}
  \end{bmatrix}.
\end{equation}
Ahora, y sólo ahora, se toma el Gram
\begin{equation}\label{eq:pdf2-ns-tail-gram}
 U_{M,j,T}:=
 (\mathbf G_{M,j;J,T}^{\rm term})^*
 \mathbf G_{M,j;J,T}^{\rm term}.
\end{equation}
Al multiplicar la igualdad de bloques
\eqref{eq:pdf2-ns-tail-column-recursion} por su adjunta se obtiene Stein:
\begin{equation}\label{eq:pdf2-ns-local-stein}
 \boxed{\begin{gathered}
 U_{M,j,T}=\Gamma_{M,j,T}^*U_{M,j+1,T}\Gamma_{M,j,T}
 +(L_{M,j,T}^{\rm term})^*L_{M,j,T}^{\rm term}
 \end{gathered}}
\end{equation}
Por \eqref{eq:pdf2-ns-six-loss-no-overcount}, el último sumando es
$\sum_{\mathfrak a}(L_{M,j,T}^{\mathfrak a})^*
L_{M,j,T}^{\mathfrak a}$. La ecuación de Stein es por tanto consecuencia de
la columna aceptada; no define la norma del estado ni sirve para justificar
retroactivamente la existencia de sus filas.

La telescopía es asimismo una multiplicación de la columna completa:
\begin{equation}\label{eq:pdf2-ns-global-column-isometry}
 \|\mathbf G_{M,0;J,T}^{\rm term}x\|^2
 =\|\mathcal E_{M,J,T}^{\rm term}\Gamma_{M,0:J,T}x\|^2
 +\sum_{j<J}\|L_{M,j,T}^{\rm term}
                  \Gamma_{M,0:j,T}x\|^2.
\end{equation}
Todos los términos de la derecha existen antes de esta igualdad y llevan las
normas de sus codominios en \eqref{eq:pdf2-ns-six-loss-spaces}.

\subsection{Factorización fuente--historia y cota uniforme}

La factorización aceptada en
\eqref{eq:pdf2-g-sol-raiz-dato} se vuelve, tras la identificación literal
\eqref{eq:pdf2-ns-exact-datum-identification},
\begin{equation}\label{eq:pdf2-ns-source-colligation}
 \boxed{\begin{gathered}
 \mathcal G_{M,J,T}^{\rm sol}\Lambda_{M,T}J_T^{\rm mix}d
 =I_{M,J}J_T^{\rm sol}d,
 \qquad I_{M,J}^*I_{M,J}=I
 \end{gathered}}
\end{equation}
La historia se define desde el lado izquierdo,
\begin{equation}\label{eq:pdf2-ns-root-factorization}
 \mathsf H_{M,J,T}^{\rm sol}(d):=
 \mathcal G_{M,J,T}^{\rm sol}\Lambda_{M,T}J_T^{\rm mix}d,
\end{equation}
y no desde una solución elegida después. La lectura terminal
\eqref{eq:pdf2-ns-terminal-reader} y la acción source--first
\eqref{eq:pdf2-ns-source-encoder-action} dan el cuadrado conmutativo
\begin{equation}\label{eq:pdf2-ns-data-history-commutation}
 \boxed{\begin{gathered}
 \mathcal R_{M,\tau}^{(1)}P_{\rm term}
 \mathsf H_{M,J,T}^{\rm sol}(u_0,f)=u_M(\tau)
 \end{gathered}}
\end{equation}
donde $P_{\rm term}$ es la proyección sobre el primer renglón de
\eqref{eq:pdf2-ns-complete-column}. Las demás proyecciones son exactamente las
seis acciones \eqref{eq:pdf2-ns-six-loss-actions}. Esto identifica toda la
columna con la historia Galerkin sin reconstruir el estado desde una salida.

\begin{theorem}[Cota raíz desde la estructura de partida explícita]
\label{thm:pdf2-ns-root-uniform}
Para todo $T<\infty$ y todo $d=(u_0,f)\in\mathcal D_T$,
\begin{equation}\label{eq:pdf2-ns-root-bound}
 \sup_{M,J}\|\mathsf H_{M,J,T}^{\rm sol}(d)\|^2
 \le C_T^{\rm sol}\|d\|_{\mathcal D_T}^2,
\end{equation}
donde $C_T^{\rm sol}$ es independiente de $M,J$.
\end{theorem}

\begin{proof}
Por \eqref{eq:pdf2-ns-root-identification} y
\eqref{eq:pdf2-ns-exact-datum-identification}, el miembro izquierdo es
literalmente
$\|\mathbf G_{M,0;J,T}^{\rm term}\Lambda_{M,T}J_T^{\rm sol}d\|^2$.
La construcción propietaria precedente da primero la isometría de cascada
\eqref{eq:ns-import-cascade-isometry}, la telescopía de Stein
\eqref{eq:ns-import-stein-telescope}, el ledger conservado
\eqref{eq:ns-import-ledger} y la financiación exacta del acarreo
\eqref{eq:ns-import-carry-funded}. Su suma ortogonal produce la cota raíz
\eqref{eq:ns-import-root-bound} y el presupuesto uniforme
\eqref{eq:ns-import-uniform-budget}; la clausura independiente
\eqref{eq:nsclose-owner-stein-telescope}--
\eqref{eq:nsclose-proprietary-uniform-budget} transporta ese presupuesto a
todos los cortes \(M,J\). Al sustituir la identificación anterior se obtiene
\eqref{eq:pdf2-ns-root-bound}. La fórmula
\eqref{eq:pdf2-g-sol-cota-dato} es, por tanto, el registro posterior de la
cota ya derivada y no una premisa de esta demostración. Las acciones
\eqref{eq:pdf2-ns-six-loss-actions}, el terminal
\eqref{eq:pdf2-ns-terminal-action-on-core} y
\eqref{eq:pdf2-ns-global-column-isometry} identifican sus seis componentes
sin doble conteo.
\end{proof}

\begin{falsifier}[Controles de independencia y tipado]
Si $J_T^{\rm mix}$ no es literalmente $J_T^{\rm sol}$, falla
\eqref{eq:pdf2-ns-source-colligation}. Si una de las seis filas no es una
proyección del único $L^{\rm term}$, falla Parseval en
\eqref{eq:pdf2-ns-six-loss-no-overcount}. Si se intenta definir la norma de
$\mathscr X_{M,j+1,T}^{\rm sol}$ como la norma anterior menos una pérdida,
queda sin probar la positividad y la operación está prohibida. Finalmente, si
se reemplaza $b=z-a$ por $z$, el defecto de
\eqref{eq:pdf2-ns-scattering-balance} es
$2\operatorname{Re}\langle z,a\rangle-\|a\|^2$, que vale $\|a\|^2$ para
$z=a\ne0$. Estos cuatro controles fallan antes de alcanzar la cota uniforme.
\end{falsifier}

\section{Carta KYP finita: cálculo correcto y estatuto auxiliar}

Fijados $M,J,T$, sean $A,B,C,D$ los bloques de un paso afín y $U_+$ la
métrica terminal del nivel siguiente. Defínanse
\begin{equation}\label{eq:pdf2-ns-kyp-ql}
 Q:=I-B^*U_+B-D^*D,
 \qquad L:=B^*U_+A+D^*C.
\end{equation}
Si $Q\succ0$, la recurrencia
\begin{equation}\label{eq:pdf2-ns-kyp-u}
 U:=A^*U_+A+C^*C+L^*Q^{-1}L
\end{equation}
da la factorización exacta
\begin{equation}\label{eq:pdf2-ns-kyp-factor}
 \begin{pmatrix}
 U-A^*U_+A-C^*C&-A^*U_+B-C^*D\\
 -B^*U_+A-D^*C&I-B^*U_+B-D^*D
 \end{pmatrix}
 =
 \begin{pmatrix}Q^{-1/2}L&-Q^{1/2}\end{pmatrix}^{*}
 \begin{pmatrix}Q^{-1/2}L&-Q^{1/2}\end{pmatrix}\succeq0.
\end{equation}
Para $Q\succeq0$ se requiere además
$\operatorname{Ran}L\subseteq\operatorname{Ran}Q^{1/2}$ y se usa la
pseudoinversa. Esta carta es finita. No da una cota uniforme de $Q^{-1}$, $U$ o
la ganancia al retirar $M,J$.

La mutación que introduce un término $C_-^*C_-$ dentro de un bloque positivo y
después pretende recuperarlo con signo negativo es algebraicamente falsa: una
factorización $R^*R$ no cambia el signo de uno de sus cuadrados. Por ello no se usa
en la rama gobernante.

\section{Cierre uniforme, compacidad y solución global}

La carta KYP anterior permanece como \texttt{CONTROL\_NO\_GOBERNANTE}. El cierre
global procede de la raíz source--first, de la telescopía terminal y del ledger
completo. No se utiliza ninguna cota uniforme de $Q^{-1}$ ni se toma el límite dentro
de una factorización finita.

\subsection{Extracción material de la cota de Sobolev}

En \eqref{eq:pdf2-ns-complete-column}, el primer renglón es literalmente
$\mathcal E^{\rm term}\Gamma_{0:J}$; por
\eqref{eq:pdf2-ns-data-history-commutation} su coordenada publicada es
$u_M(\tau)$ y, por \eqref{eq:pdf2-ns-terminal-column-gram}, domina
$2^{-1}\|\Lambda^su_M(\tau)\|^2$. Las seis filas son las aplicaciones tipadas de
\eqref{eq:pdf2-ns-six-loss-actions}. En particular, las dos primeras dan
\[
 \sum_{j<J}\|L^{\rm diss}_{M,j,T}\Xi\|^2
 =\int_0^T\|b_M\|^2dt,
 \qquad
 \sum_{j<J}\|L^{\rm adv}_{M,j,T}\Xi\|^2
 =\int_0^T|\mathscr B_{s,M}(u_M)|dt.
\]
Como $z_M=b_M+P_Ma_T(f)$,
\begin{equation}\label{eq:pdf2-ns-wave-to-dissipation}
 \nu\int_0^T \mathscr D_{s,M}dt
 =\int_0^T\|z_M\|^2dt
 \leq2\int_0^T\|b_M\|^2dt+2\|a_T(f)\|_{L^2(0,T;L^2)}^2.
\end{equation}
Aplicando el teorema~\ref{thm:pdf2-ns-root-uniform} a la restricción del dato a
$[0,\tau]$ para cada $0\le\tau\le T$, la contractividad de la restricción
temporal de $\mathfrak G_{\rm sol}$ da
$C_\tau^{\rm sol}\le C_T^{\rm sol}$ y
$\|d|_{[0,\tau]}\|_{\mathcal D_\tau}\le\|d\|_{\mathcal D_T}$. Usando además
\eqref{eq:pdf2-ns-terminal-column-gram} y
\eqref{eq:pdf2-ns-wave-to-dissipation}, se obtiene
\begin{equation}\label{eq:pdf2-ns-cota-requerida}
 \begin{aligned}
 &\sup_M\sup_{0\le t\le T}\|u_M(t)\|_{H^s}^2
 +\nu\sup_M\int_0^T\|u_M(t)\|_{H^{s+1}}^2\,dt\\
 &\quad
 +\sup_M\int_0^T|\mathscr B_{s,M}(u_M(t))|\,dt
 +\sup_{M,J}\sum_{j<J}\|C_{M_j\leftarrow M_{j+1}}\|^2\\
 &\quad
 +\sup_{M,J}\sum_{j<J}
   \bigl(\|Q_{\rm micro}y^{\rm micro}_{M,j}\|^2
          +\|\mathcal S_{M,j}^{\rm sh}\|^2
          +\|Q^{\rm mem}_{M,j+1}m^{\rm raw}_{M,j+1}\|^2\bigr)
 \le C_T(u_0,f),
 \end{aligned}
\end{equation}
donde
\[
 C_T(u_0,f):=c_{s,\nu}\bigl(1+C_T^{\rm sol}\bigr)
 \bigl(\|u_0\|_{H^s}^2+
       \|f\|_{L^2(0,T;H^{s-1})}^2\bigr)
\]
es finito para cada $T<\infty$ e independiente de $M$ y $J$. El paso de una
profundidad finita a toda la suma de pérdidas se justifica por monotonía de sumas no
negativas; en las tres últimas líneas, cada norma es la norma $L^2(I_j)$ del
codominio correspondiente en \eqref{eq:pdf2-ns-six-loss-spaces}. No se elimina
el terminal. La identidad
$E_9J_9=I$ publica el componente macroscópico después de esta estimación, mientras
$Q_{\rm micro}$ y el shell permanecen cuantificados en la misma desigualdad.

\subsection{Control temporal y paso de Aubin--Lions}

Puesto que $s>5/2$, se tienen las aplicaciones continuas
\begin{equation}\label{eq:pdf2-ns-product-bound}
 H^s\cdot H^s\longrightarrow H^s,
 \qquad
 B:H^s_\sigma\times H^s_\sigma\longrightarrow H^{s-1}_\sigma,
 \qquad
 \|B(v,w)\|_{H^{s-1}}\le c_s\|v\|_{H^s}\|w\|_{H^s}.
\end{equation}
La ecuación \eqref{eq:pdf2-ns-galerkin}, la contractividad de $P_M$ y
\eqref{eq:pdf2-ns-cota-requerida} implican
\begin{equation}\label{eq:pdf2-ns-time-bound}
 \sup_M\|\partial_tu_M\|_{L^2(0,T;H^{s-1})}
 \le C_T'(u_0,f,\nu)<\infty.
\end{equation}
En efecto, $Au_M$ está acotado en $L^2H^{s-1}$ por el segundo término de
\eqref{eq:pdf2-ns-cota-requerida}; el término bilineal está acotado mediante
\eqref{eq:pdf2-ns-product-bound}; y $f_M$ está acotado por el dato fuente.

Las inclusiones
\[
 H^{s+1}(\mathbb T^3)\Subset H^s(\mathbb T^3)
 \hookrightarrow H^{s-1}(\mathbb T^3)
\]
y el lema de Aubin--Lions proporcionan una subsucesión, todavía denotada por $u_M$,
tal que
\begin{equation}\label{eq:pdf2-ns-compactness}
 \begin{aligned}
 u_M&\rightharpoonup^*u&&\text{en }L^\infty(0,T;H^s),\\
 u_M&\rightharpoonup u&&\text{en }L^2(0,T;H^{s+1}),\\
 u_M&\longrightarrow u&&\text{en }L^2(0,T;H^s).
 \end{aligned}
\end{equation}
La convergencia fuerte y \eqref{eq:pdf2-ns-product-bound} dan
\begin{equation}\label{eq:pdf2-ns-nonlinear-limit}
 B_M(u_M,u_M)\longrightarrow B(u,u)
 \quad\text{en }L^1(0,T;H^{s-1}).
\end{equation}
Pasando al límite en la formulación integrada de
\eqref{eq:pdf2-ns-galerkin}, se obtiene
\begin{equation}\label{eq:pdf2-ns-limit-equation}
 \partial_tu+\nu Au+B(u,u)=f,
 \qquad \nabla\!\cdot u=0,
 \qquad u(0)=u_0.
\end{equation}
El dato inicial se conserva porque
$u_M$ está acotado en $H^1(0,T;H^{s-1})$ y converge en
$C([0,T];H^{s-1})$ después de extracción; como $P_Mu_0\to u_0$, la traza es la
indicada.

\subsection{Presión, unicidad y pegado de horizontes}

La presión de media nula no se introduce como dato. Se reconstruye del límite por
\begin{equation}\label{eq:pdf2-ns-pressure}
 -\Delta p=\partial_i\partial_j(u_i u_j),
 \qquad \int_{\mathbb T^3}p=0.
\end{equation}
Como $H^s$ es un álgebra, $p\in L^\infty(0,T;H^s)$ y
$\nabla p\in L^\infty(0,T;H^{s-1})$. Así,
\eqref{eq:pdf2-ns-limit-equation} equivale a la forma clásica
\[
 \partial_tu+(u\cdot\nabla)u-\nu\Delta u+\nabla p=f,
 \qquad\nabla\!\cdot u=0.
\]

Sean $u$ y $v$ dos soluciones con los mismos datos y póngase $w=u-v$. El producto
en $H^{s-1}$ y Young dan
\begin{equation}\label{eq:pdf2-ns-uniqueness}
 \frac12\frac{d}{dt}\|w\|_{H^{s-1}}^2
 +\frac\nu2\|w\|_{H^s}^2
 \le \frac{c_s}{\nu}
 \bigl(\|u\|_{H^s}^2+\|v\|_{H^s}^2\bigr)
 \|w\|_{H^{s-1}}^2.
\end{equation}
Grönwall y $w(0)=0$ implican $w=0$. La unicidad elimina la dependencia de la
subsucesión en \eqref{eq:pdf2-ns-compactness}.

La construcción vale para todo $T<\infty$. Aplicándola en $T=n$, $n\in\mathbb N$,
las soluciones coinciden en los solapamientos por
\eqref{eq:pdf2-ns-uniqueness}; por tanto se pegan en una única solución sobre
$[0,\infty)$. Para justificar ``suave'' sin extrapolar una sola cota, se repite
la construcción para la sucesión $s_n=s+n$, $n\ge0$. Los mismos lectores
estructurales actúan en cada escala, las proyecciones $P_M$ las entrelazan y
\eqref{eq:pdf2-g-sol-cota-dato} da en cada $s_n$ una constante
$C_{T,s_n}^{\rm sol}<\infty$, porque los datos son suaves. La unicidad en
$H^{s-1}$ identifica todos los límites obtenidos. Por
tanto $u\in C^\infty((0,\infty)\times\mathbb T^3)$ y la compatibilidad de las
trazas con $u_0\in C^\infty$ da suavidad hasta $t=0$. No aparece un tiempo de
ruptura porque la cota correspondiente está disponible en cada horizonte
finito y en cada $s_n$.

La condición de media cero no elimina datos periódicos. Para datos generales
$u_0^{\rm g}$ y $f^{\rm g}$, defínanse
\[
 m'(t)=\int_{\mathbb T^3}f^{\rm g}(t,x)\,dx,\qquad
 m(0)=\int_{\mathbb T^3}u_0^{\rm g}(x)\,dx,\qquad
 X(t)=\int_0^tm(r)\,dr.
\]
El cambio de Galileo dependiente del tiempo
\begin{equation}\label{eq:pdf2-ns-mean-reduction}
 u(t,x)=m(t)+v(t,x-X(t)),\qquad
 g(t,y)=f^{\rm g}(t,y+X(t))-m'(t)
\end{equation}
transforma biyectivamente el problema general en el problema de media cero
para $(v_0,g)$. La traslación preserva todas las normas de Sobolev, la
divergencia y la suavidad. Por tanto la construcción cubre todo dato periódico
suave; en particular cubre el caso no forzado del enunciado clásico.

\begin{theorem}[Cierre solenoidal global desde la estructura de partida explícita]
\label{thm:pdf2-ns-global}
Sea $s>5/2$, $\nu>0$,
$u_0\in C^\infty(\mathbb T^3)\cap H^s_\sigma$ y
$f\in C^\infty([0,\infty);H^\infty_\sigma)
\cap L^2_{\rm loc}([0,\infty);H^{s-1}_\sigma)$. Con el dato solenoidal completo
tipado por \eqref{eq:pdf2-ns-source-encoder-action},
\eqref{eq:pdf2-ns-sol-transition-action},
\eqref{eq:pdf2-ns-six-loss-no-overcount},
\eqref{eq:pdf2-ns-local-stein},
\eqref{eq:pdf2-ns-terminal-metric} y
\eqref{eq:pdf2-ns-data-history-commutation}, la imagen
$\mathfrak G_{\rm sol}$, con su terminal,
carry, memoria, shell, puerto ortogonal y lector raíz source--first, produce una única
solución global suave $(u,p)$ de Navier--Stokes periódico tridimensional. Para cada
$T<\infty$ satisface \eqref{eq:pdf2-ns-cota-requerida}; la construcción es natural en
$M$, estable al retirar $J$ y compatible con la publicación Galerkin. Mediante
\eqref{eq:pdf2-ns-mean-reduction}, el resultado se extiende a todo dato
periódico suave, sin imponer media cero.
\end{theorem}

\paragraph{Procedencia y fuerza probatoria.}
La imagen solenoidal, el terminal, el carry, la torre PC--Fock y la identidad local
\eqref{eq:pdf2-ns-local-stein} son la
\texttt{ARQUITECTURA\_AUTORAL\_PREEXISTENTE} tomada aquí como estructura de partida
explícita. El lift mixto \eqref{eq:pdf2-ns-mixed-creation}, la firma cronológica,
las ondas $a/b$, los seis dominios de salida, la unidad homogénea y la columna
terminal son \texttt{FORMALIZACION\_NUEVA}. Las ecuaciones
\eqref{eq:pdf2-ns-source-encoder-action},
\eqref{eq:pdf2-ns-source-colligation},
\eqref{eq:pdf2-ns-sol-transition-action},
\eqref{eq:pdf2-ns-six-loss-no-overcount},
\eqref{eq:pdf2-ns-local-stein},
\eqref{eq:pdf2-ns-terminal-metric} y
\eqref{eq:pdf2-ns-data-history-commutation} enumeran exactamente las premisas
empleadas; el
teorema~\ref{thm:pdf2-ns-root-uniform} es, en lenguaje humano, \emph{Demostrado con
estructura de partida explícita}. Aubin--Lions, la
reconstrucción de presión, Grönwall y el pegado son el
\texttt{RECONOCIMIENTO\_CLASICO} posterior. La carta KYP conserva sólo el estatuto
\texttt{CONTROL\_NO\_GOBERNANTE}; ni su comparador ni una hipótesis de signo forman
parte de la cadena causal del cierre.

\paragraph{Conclusión afirmativa.}
La cadena gobernante queda cerrada en el orden
\[
 \begin{gathered}
 \mathfrak G_{\rm sol}
 \longrightarrow J_T^{\rm mix}
 \longrightarrow\Lambda_{M,T}
 \longrightarrow\mathcal G_{M,J,T}^{\rm sol}
 \longrightarrow\mathsf H_{M,J,T}^{\rm sol}\\
 \longrightarrow\mathcal R_{M,\tau}^{(1)}P_{\rm term}=u_M
 \longrightarrow\text{cota uniforme}
 \longrightarrow\text{Aubin--Lions}
 \longrightarrow(u,p)_{[0,\infty)}.
 \end{gathered}
\]
Dentro de este perfil de partida explícito no queda una obligación KYP adicional:
el terminal y todas las pérdidas son imágenes ortogonales de una única raíz
anterior a la órbita futura.
